% GENERATED FILE. Edit canonical lessons, then run npm run generate:books.
% canonical-source-hash: fnv1a64:cca640f1c29b4584
% canonical-lessons: MW-C237-titar, MW-C237-pakh, MW-C237-chugno, MW-C237-indo, MW-C237-midko

\chapter{Small Creatures of the Field}
\label{ch:mw-small-creatures-of-the-field}
\hlchaptermodality{237}

\begin{chapteropening}
\textbf{By the end of this chapter:} \emph{I can name a partridge, a frog, wings and eggs.}

Complete the last lesson of chapter 237: I can name a partridge, a frog, wings and eggs.
\end{chapteropening}

\section[\texorpdfstring{\mw{तीतर} (tītar)}{तीतर (tītar)}]{\mw{तीतर} (tītar) — a partridge}

\label{lesson:MW-C237-titar}



\begin{quote}
Before the new one: say the Marwadi for a horn, then the Marwadi for a hoof.
\end{quote}

\subsection*{You'll want to know: \mw{तीतर}}
\textbf{\mw{तीतर}} — \emph{tītar} — \textquotedblleft{}a partridge\textquotedblright{}.

\begin{quote}
\mw{खेत} \mw{मांय} \mw{तीतर} \mw{है।}
\end{quote}

\begin{quote}
\emph{khet mā̃y tītar hai.}
\end{quote}

\begin{quote}
There's a partridge in the field.
\end{quote}

A masculine noun.

\begin{cousinweb}
\textbf{\mw{तीतर}} (\emph{tītar}) is inherited from Sanskrit \emph{tittira}, \textquotedblleft{}a partridge\textquotedblright{} (CDIAL 5809, which lists Marwari \emph{tītar}).
\end{cousinweb}

\begin{grammarlens}[title={What you've built}]
One more word for this chapter.
\end{grammarlens}

\subsection*{Your turn}
\begin{itemize}
\raggedright
  \item \emph{Say it:} \emph{tītar}
  \item \emph{Say it:} \emph{tītar}, once more
  \item \emph{Say it:} \emph{khet mā̃y tītar hai.} — There's a partridge in the field.
  \item \emph{Recall:} say the Marwadi for a horn, then the Marwadi for a hoof, then say \emph{tītar} again
\end{itemize}

\begin{tcolorbox}[breakable,colback=teal!4,colframe=teal!35!black,title={Before you move on}]
What does \mw{तीतर} mean? (\textquotedblleft{}A partridge\textquotedblright{}.) Say it once more.
\end{tcolorbox}
\section[\texorpdfstring{\mw{पांख} (pā̃kh)}{पांख (pā̃kh)}]{\mw{पांख} (pā̃kh) — a wing; a feather}

\label{lesson:MW-C237-pakh}



\begin{quote}
Before the new one: say the Marwadi for a hoof, then the Marwadi for a partridge.
\end{quote}

\subsection*{You'll want to know: \mw{पांख}}
\textbf{\mw{पांख}} — \emph{pā̃kh} — \textquotedblleft{}a wing; a feather\textquotedblright{}.

\begin{quote}
\mw{चिड़ी} \mw{पांख} \mw{खोलै} \mw{है।}
\end{quote}

\begin{quote}
\emph{chiṛī pā̃kh kholai hai.}
\end{quote}

\begin{quote}
The sparrow spreads its wings.
\end{quote}

A feminine noun.

\begin{cousinweb}
\textbf{\mw{पांख}} (\emph{pā̃kh}) is inherited from Sanskrit \emph{pakṣa}, \textquotedblleft{}a wing\textquotedblright{}; CDIAL lists an Old Marwari word for a wing or a feather under it (CDIAL 7627).
\end{cousinweb}

\begin{grammarlens}[title={What you've built}]
One more word for this chapter.
\end{grammarlens}

\subsection*{Your turn}
\begin{itemize}
\raggedright
  \item \emph{Say it:} \emph{pā̃kh}
  \item \emph{Say it:} \emph{pā̃kh}, once more
  \item \emph{Say it:} \emph{chiṛī pā̃kh kholai hai.} — The sparrow spreads its wings.
  \item \emph{Recall:} say the Marwadi for a hoof, then the Marwadi for a partridge, then say \emph{pā̃kh} again
\end{itemize}

\begin{tcolorbox}[breakable,colback=teal!4,colframe=teal!35!black,title={Before you move on}]
What does \mw{पांख} mean? (\textquotedblleft{}A wing; a feather\textquotedblright{}.) Say it once more.
\end{tcolorbox}
\section[\texorpdfstring{\mw{चुगणो} (chugṇo)}{चुगणो (chugṇo)}]{\mw{चुगणो} (chugṇo) — to peck; to pick up (grain or berries)}

\label{lesson:MW-C237-chugno}



\begin{quote}
Before the new one: say the Marwadi for a partridge, then the Marwadi for a wing.
\end{quote}

\subsection*{You'll want to know: \mw{चुगणो}}
\textbf{\mw{चुगणो}} — \emph{chugṇo} — \textquotedblleft{}to peck; to pick up (grain or berries)\textquotedblright{}.

\begin{quote}
\mw{कबूतर} \mw{बाजरो} \mw{चुगै} \mw{है।}
\end{quote}

\begin{quote}
\emph{kabūtar bājro chugai hai.}
\end{quote}

\begin{quote}
The pigeon pecks at the millet.
\end{quote}

Past (I …): \mw{म्हैं} \mw{केर} \mw{चुग्या}, \emph{mhaĩ ker chugyā}.

\begin{cousinweb}
\textbf{\mw{चुगणो}} (\emph{chugṇo}) is inherited from a reconstructed \emph{*cugyati}, \textquotedblleft{}pecks\textquotedblright{}; CDIAL lists Marwari \emph{cugṇo}, \textquotedblleft{}to peck, to eat\textquotedblright{} (CDIAL 4852).
\end{cousinweb}

\begin{grammarlens}[title={What you've built}]
One more word for this chapter.
\end{grammarlens}

\subsection*{Your turn}
\begin{itemize}
\raggedright
  \item \emph{Say it:} \emph{chugṇo}
  \item \emph{Say it:} \emph{chugṇo}, once more
  \item \emph{Say it:} \emph{kabūtar bājro chugai hai.} — The pigeon pecks at the millet.
  \item \emph{Recall:} say the Marwadi for a partridge, then the Marwadi for a wing, then say \emph{chugṇo} again
\end{itemize}

\begin{tcolorbox}[breakable,colback=teal!4,colframe=teal!35!black,title={Before you move on}]
What does \mw{चुगणो} mean? (\textquotedblleft{}To peck; to pick up (grain or berries)\textquotedblright{}.) Say it once more.
\end{tcolorbox}
\section[\texorpdfstring{\mw{इंडो} (iṇḍo)}{इंडो (iṇḍo)}]{\mw{इंडो} (iṇḍo) — an egg}

\label{lesson:MW-C237-indo}



\begin{quote}
Before the new one: say the Marwadi for a wing, then the Marwadi for to peck.
\end{quote}

\subsection*{You'll want to know: \mw{इंडो}}
\textbf{\mw{इंडो}} — \emph{iṇḍo} — \textquotedblleft{}an egg\textquotedblright{}.

\begin{quote}
\mw{चिड़ी} \mw{इंडो} \mw{देवै} \mw{है।}
\end{quote}

\begin{quote}
\emph{chiṛī iṇḍo devai hai.}
\end{quote}

\begin{quote}
The sparrow lays an egg.
\end{quote}

A masculine noun.

\begin{cousinweb}
\textbf{\mw{इंडो}} (\emph{iṇḍo}) is inherited from Sanskrit \emph{aṇḍa}, \textquotedblleft{}an egg\textquotedblright{}; CDIAL lists Gujarati \emph{ĩḍũ}, \textquotedblleft{}an egg\textquotedblright{}. CDIAL gives no Marwari form in this entry (CDIAL 1111).
\end{cousinweb}

\begin{grammarlens}[title={What you've built}]
One more word for this chapter.
\end{grammarlens}

\subsection*{Your turn}
\begin{itemize}
\raggedright
  \item \emph{Say it:} \emph{iṇḍo}
  \item \emph{Say it:} \emph{iṇḍo}, once more
  \item \emph{Say it:} \emph{chiṛī iṇḍo devai hai.} — The sparrow lays an egg.
  \item \emph{Recall:} say the Marwadi for a wing, then the Marwadi for to peck, then say \emph{iṇḍo} again
\end{itemize}

\begin{tcolorbox}[breakable,colback=teal!4,colframe=teal!35!black,title={Before you move on}]
What does \mw{इंडो} mean? (\textquotedblleft{}An egg\textquotedblright{}.) Say it once more.
\end{tcolorbox}
\section[\texorpdfstring{\mw{मींडको} (mī̃ḍko)}{मींडको (mī̃ḍko)}]{\mw{मींडको} (mī̃ḍko) — a frog}

\label{lesson:MW-C237-midko}



\begin{quote}
Before the new one: say the Marwadi for to peck, then the Marwadi for an egg.
\end{quote}

\subsection*{You'll want to know: \mw{मींडको}}
\textbf{\mw{मींडको}} — \emph{mī̃ḍko} — \textquotedblleft{}a frog\textquotedblright{}.

\begin{quote}
\mw{बरसात} \mw{मांय} \mw{मींडको} \mw{बोलै} \mw{है।}
\end{quote}

\begin{quote}
\emph{barsāt mā̃y mī̃ḍko bolai hai.}
\end{quote}

\begin{quote}
In the rain the frog croaks.
\end{quote}

A masculine noun.

\begin{cousinweb}
\textbf{\mw{मींडको}} (\emph{mī̃ḍko}) is inherited from Sanskrit \emph{maṇḍūka}, \textquotedblleft{}a frog\textquotedblright{}; CDIAL lists Old Marwari \emph{mīḍako}, \textquotedblleft{}a frog\textquotedblright{} (CDIAL 9746).
\end{cousinweb}

\begin{grammarlens}[title={What you've built}]
One more word for this chapter.
\end{grammarlens}

\subsection*{Your turn}
\begin{itemize}
\raggedright
  \item \emph{Say it:} \emph{mī̃ḍko}
  \item \emph{Say it:} \emph{mī̃ḍko}, once more
  \item \emph{Say it:} \emph{barsāt mā̃y mī̃ḍko bolai hai.} — In the rain the frog croaks.
  \item \emph{Recall:} say the Marwadi for to peck, then the Marwadi for an egg, then say \emph{mī̃ḍko} again
\end{itemize}

\begin{tcolorbox}[breakable,colback=teal!4,colframe=teal!35!black,title={Before you move on}]
What does \mw{मींडको} mean? (\textquotedblleft{}A frog\textquotedblright{}.) Say it once more.
\end{tcolorbox}
